\section{Lattice Yang--Mills Theory}
\label{sec:lattice}
\label{sec:fundamental-proof}
%=============================================================================

This section provides the rigorous mathematical construction of lattice Yang--Mills 
theory. All definitions are precise and all statements are either definitions, 
standard mathematical facts, or proven theorems.

\subsection{The Lattice}

\begin{definition}[Hypercubic Lattice]
\label{def:lattice}
For $L \in \mathbb{N}$, the finite hypercubic lattice is:
\[
\Lambda_L = (\mathbb{Z}/L\mathbb{Z})^4
\]
with periodic boundary conditions. The lattice has:
\begin{itemize}
    \item Vertices: $V_L = \Lambda_L$ with $|V_L| = L^4$
    \item Oriented edges: $E_L$ with $|E_L| = 4L^4$
    \item Plaquettes (unit squares): $P_L$ with $|P_L| = 6L^4$
\end{itemize}
\end{definition}

\begin{definition}[Edges and Orientation]
An edge $e = (x, \mu)$ connects vertex $x$ to vertex $x + \hat{\mu}$, where 
$\hat{\mu}$ is the unit vector in direction $\mu \in \{1,2,3,4\}$. 
The reverse edge is $\bar{e} = (x + \hat{\mu}, -\mu)$.
\end{definition}

\begin{definition}[Plaquette]
A plaquette $p = (x, \mu, \nu)$ with $\mu < \nu$ is the unit square with corners 
$x$, $x + \hat{\mu}$, $x + \hat{\mu} + \hat{\nu}$, $x + \hat{\nu}$.
\end{definition}

\subsection{Gauge Fields}

\begin{definition}[Configuration Space]
\label{def:config-space}
A gauge field configuration is a map:
\[
U : E_L \to SU(N)
\]
satisfying $U_{\bar{e}} = U_e^{-1}$ (orientation reversal).

The configuration space is:
\[
\mathcal{A}_L = SU(N)^{|E_L|/2} \cong SU(N)^{2L^4}
\]
which is a compact manifold of dimension $2L^4(N^2-1)$.
\end{definition}

\begin{definition}[Plaquette Variable]
\label{def:plaquette-var}
For plaquette $p = (x, \mu, \nu)$, the plaquette variable (holonomy) is:
\[
U_p = U_{(x,\mu)} U_{(x+\hat{\mu},\nu)} U_{(x+\hat{\mu}+\hat{\nu},-\mu)} U_{(x+\hat{\nu},-\nu)}
\]
This is the ordered product of link variables around the plaquette boundary.
\end{definition}

\begin{lemma}[Plaquette Properties]
\label{lem:plaquette-props}
For any configuration $U$:
\begin{enumerate}[label=(\alph*)]
    \item $U_p \in SU(N)$
    \item Under gauge transformation $g : V_L \to SU(N)$:
    \[
    U_p \mapsto g(x) U_p g(x)^{-1}
    \]
    \item $\Tr(U_p)$ is gauge-invariant
\end{enumerate}
\end{lemma}

\begin{proof}
(a) follows from closure of $SU(N)$ under multiplication.
(b) follows by direct computation using $U_e \mapsto g(x)U_e g(y)^{-1}$ for $e = (x,y)$.
(c) follows from (b) and cyclicity of trace.
\end{proof}

\subsection{The Wilson Action}

\begin{definition}[Wilson Action]
\label{def:wilson-action}
\label{def:wilson_action}
The Wilson action is:
\[
S_W[U] = \frac{\beta}{N} \sum_{p \in P_L} \Re \Tr(U_p)
\]
where $\beta > 0$ is the coupling constant and the sum is over all plaquettes.
\end{definition}

\begin{lemma}[Action Properties]
\label{lem:action-props}
The Wilson action satisfies:
\begin{enumerate}[label=(\alph*)]
    \item \textbf{Gauge Invariance:} $S_W[U^g] = S_W[U]$ for all gauge transformations $g$.
    \item \textbf{Boundedness:} $-6L^4\beta \leq S_W[U] \leq 6L^4\beta$.
    \item \textbf{Maximum:} $S_W[U] = 6L^4\beta$ if and only if $U_p = \mathbf{1}$ for all $p$.
    \item \textbf{Continuity:} $S_W : \mathcal{A}_L \to \mathbb{R}$ is continuous.
\end{enumerate}
\end{lemma}

\begin{proof}
(a) follows from Lemma~\ref{lem:plaquette-props}(c).
(b) follows from $|\Re\Tr(U)| \leq N$ for $U \in SU(N)$.
(c) follows from $\Re\Tr(U) = N$ iff $U = \mathbf{1}$.
(d) follows from continuity of trace and multiplication on $SU(N)$.
\end{proof}

\subsection{The Partition Function and Measure}

\begin{definition}[Haar Measure]
\label{def:haar}
Let $dU$ denote the normalized Haar measure on $SU(N)$, satisfying:
\begin{enumerate}[label=(\alph*)]
    \item $\int_{SU(N)} dU = 1$ (normalization)
    \item $\int f(gU) \, dU = \int f(Ug) \, dU = \int f(U) \, dU$ (bi-invariance)
\end{enumerate}
\end{definition}

\begin{definition}[Product Measure]
\label{def:product-measure}
The reference measure on $\mathcal{A}_L$ is the product Haar measure:
\[
\mathcal{D}U = \prod_{e \in E_L^+} dU_e
\]
where $E_L^+ \subset E_L$ contains one edge from each pair $(e, \bar{e})$.
\end{definition}

\begin{definition}[Partition Function]
\label{def:partition}
The partition function is:
\[
Z_L(\beta) = \int_{\mathcal{A}_L} e^{S_W[U]} \, \mathcal{D}U
\]
\end{definition}

\begin{proposition}[Partition Function Properties]
\label{prop:partition-props}
The partition function satisfies:
\begin{enumerate}[label=(\alph*)]
    \item $Z_L(\beta) > 0$ for all $\beta > 0$ and $L \geq 1$.
    \item $Z_L(\beta)$ is real-analytic in $\beta$ on $(0, \infty)$.
    \item $Z_L(0) = 1$ (product of normalized Haar measures).
    \item $e^{-6L^4\beta} \leq Z_L(\beta) \leq e^{6L^4\beta}$.
\end{enumerate}
\end{proposition}

\begin{proof}
(a) The integrand $e^{S_W} > 0$ and $\mathcal{D}U$ is a probability measure on a compact space.
(b) The integrand is real-analytic in $\beta$ uniformly on the compact domain.
(c) At $\beta = 0$: $e^{S_W} = 1$, so $Z_L(0) = \int \mathcal{D}U = 1$.
(d) Follows from Lemma~\ref{lem:action-props}(b).
\end{proof}

\begin{definition}[Gibbs Measure]
\label{def:gibbs}
The finite-volume Gibbs measure is:
\[
\langle \cdot \rangle_L = \frac{1}{Z_L(\beta)} \int (\cdot) \, e^{S_W[U]} \, \mathcal{D}U
\]
\end{definition}

\subsection{Wilson Loops and Observables}

\begin{definition}[Wilson Loop]
\label{def:wilson-loop}
For a closed curve $\gamma$ on the lattice (sequence of edges forming a loop), 
the Wilson loop is:
\[
W_\gamma[U] = \Tr\left( \prod_{e \in \gamma} U_e \right)
\]
where the product respects orientation.
\end{definition}

\begin{definition}[Rectangular Wilson Loop]
\label{def:rectangular-loop}
The rectangular Wilson loop $W_{R,T}$ is the trace of the holonomy around a 
rectangle of spatial extent $R$ and temporal extent $T$.
\end{definition}

\begin{definition}[String Tension]
\label{def:string-tension-sec02}
The string tension is defined as:
\[
\sigma(\beta) = -\lim_{R,T \to \infty} \frac{1}{RT} \log \langle W_{R,T} \rangle
\]
when this limit exists.
\end{definition}

\begin{remark}
The existence of this limit and the area law $\langle W_{R,T} \rangle \sim e^{-\sigma RT}$ 
is a statement about \emph{confinement}. It is proven for $\beta < \beta_0$ 
(Theorem~\ref{thm:strong-coupling-string}) but remains open for general $\beta$.
\end{remark}

\subsection{Gauge Invariance}
\label{sec:gauge-invariance}

\begin{definition}[Gauge Transformation]
\label{def:gauge-transf}
A gauge transformation is a map $g : V_L \to SU(N)$. It acts on configurations by:
\[
(U^g)_e = g(x) U_e g(y)^{-1}
\]
for edge $e = (x, y)$.
\end{definition}

\begin{proposition}[Gauge Invariance of Measure]
\label{prop:gauge-inv-measure}
The Gibbs measure is gauge-invariant:
\[
\langle f(U^g) \rangle = \langle f(U) \rangle
\]
for any gauge transformation $g$ and any measurable $f$.
\end{proposition}

\begin{proof}
The Haar measure $dU_e$ is bi-invariant, so $\mathcal{D}U$ is invariant under 
$U \mapsto U^g$. The Wilson action is gauge-invariant (Lemma~\ref{lem:action-props}(a)).
\end{proof}

\subsection{Free Energy}

\begin{definition}[Free Energy Density]
\label{def:free-energy}
The free energy density is:
\[
f_L(\beta) = -\frac{1}{|P_L|} \log Z_L(\beta) = -\frac{1}{6L^4} \log Z_L(\beta)
\]
\end{definition}

\begin{proposition}[Bounds on Free Energy]
For all $\beta > 0$ and $L \geq 1$:
\[
-\beta \leq f_L(\beta) \leq \beta
\]
\end{proposition}

\begin{proof}
Immediate from Proposition~\ref{prop:partition-props}(d).
\end{proof}

\begin{conjecture}[Thermodynamic Limit]
\label{conj:thermo-limit}
The limit
\[
f(\beta) = \lim_{L \to \infty} f_L(\beta)
\]
exists for all $\beta > 0$.
\end{conjecture}

\textbf{Status:} Proven for $\beta < \beta_0$ (strong coupling). \textbf{Open} for 
general $\beta$.


